% Pursuing Stacks (cotes 134-2 to 134-8): archive cote 134-2, site page 80 (PDF page 81).
% notes to ch. I. Draft compared: ps1.tex:1188-1193 ; ps1.tex:1193-1195 ; ps1.tex:1205-1205.
% AI-transcribed by Claude Opus 5.5 (claude-opus-5-5, adaptive thinking, effort high), Message Batches API, batch ps1 (msgbatch_01Rv5hCaNTtUt4JHx1yqoHPL), 2026-11-12; not yet checked by eye.
\scanpage{134-2/80}
\mspage{64}

(${}^{22}$) For the notion of a ``\uline{lien}'' \hand{(}or ``tie''\hand{)}\hand{,} wh\over{a}{\unsure{i}}ch is one of the main in-%
gredients of the non commutative cohomology panoply of Giraud's theory, I
refer to his book (Springer\hand{,}Grundlehren 179, 1971).\hand{/}\inserted{A} \uline{Picard category}\xout{xx}
\margin{$\otimes$,}\inserted[typed]{is a groupoi} \xout{\unsure{is a category is}} endowed with an operation \inserted{$\otimes$} \deleted{\illegible{1}} \xout{which makes it resemble to an}
together with associativity, unity and commutativity data for this opera-%
tion, which make it resemble to a commutative group. \inserted{A} ``\uline{Champ\xout{x} de Picard}''
\hand{(}or ``Picard stack\deleted{\unsure{s}}''\hand{)} \inserted[typed]{is} \xout{are} defined accordingly, by relativizing over an arbi-%
trary space or topos (replacing the groupoid by a stack of groupoids over
this topos). The necessary \hand{``}general\deleted{\unsure{ly}} nonsense\hand{''} on these \inserted{notions} is \unsure{developped|developed} rathe
carefully in an exposé of Deligne in SGA 4 (SGA 4 XVIII 1.4). In this
letter to Larry Breen, I am assuming ``known'' the notion of an n-stack
(for $n$\over{$=2$}{\unsure{$\leq 3$}} at any rate), and the corresponding notion of (strict) \deleted{n-}\uline{Picard}
\inserted{n-}\uline{stack}, which should be describable (as was explained in Deligne's notes
in the case n=1) by an n-truncated chain complex in the category of abeli-%
an sheaves on X (viewed mainly as an object of the relevant derived catego-%
ry). The ``strictness'' condition on usual Picard stacks refers to the
restriction that the commutativity isomorphism within an object $L$\hand{$\otimes$}$L'$ , \xout{wh}
when L=L', should reduce to the identity. It is assumed (without further
explanation) that the condition carries over in a natural way to Picard
n-stacks, in such a way as to allow an interpretation of these by truncate\hand{d}
objects in a suitable derived category, as hinted above.

(${}^{23}$) When M is any abelian sheaf on a topos, \inserted{(the ``Mac-Lane resolution'')} X(M) is a certain canoni-%
cal \inserted{left} resolution of M by sheaves of $\underline{Z}$-modules which are ``free'', and more
specifically, which are finite direct sums of sheaves of the type
$\underline{Z}^{(T)}$ , where T is any sheaf of the type $M^{n}$ (finite product of copies of
M ). This canonical\deleted{ly} construction was introduced \deleted{first} by MacLane\inserted{(for abelian groups)}, and
gained new popularity in the French school of algebraic geometry and homo-%
logical algebra in the late sixties, because \xout{in case M is described by}
\xout{some abelian group scheme over a ground} it gives a very handy way to rela-%
te the $\mathrm{Ext}^{i}(M,N)$ invariants (when N is another abelian sheaf on X) to
the \xout{\unsure{sp}} ``spacial'' cohomology of M (i.e. of the induced topos $X_{/M}$\hand{)} with
coefficients in N.

(${}^{24}$) Reflecting on the ``right'' version of the provisional Verdier notion
of a \inserted[typed]{triangulated} \xout{derived} category (which was supposed to describe adequately\deleted{\unsure{,}} the \inserted{relevant} inter-%
nal structure of the derived categories of abelian categories) is part of
my present program for the notes on Pursuing Stacks, and will be the main
task in one of the chapters of volume two. \deleted{Compare} \inserted{For} \hand{some indications along}
\hand{these lines, see also section \unsure{69} (sketching the basic notion of a ``\uline{derivator}'').}

(${}^{25}$) As was seen in section 9, ``uniqueness'' here has to be understood in

\margin{80}
